\documentclass[11pt]{article}

\usepackage[margin=1in]{geometry}
\usepackage{amsmath}
\usepackage{amssymb}
\usepackage{amsthm}
\usepackage{mathtools}
\usepackage{xcolor}
\usepackage[colorlinks=true,linkcolor=blue,urlcolor=blue]{hyperref}
\usepackage[final]{microtype}

\title{heterosplat: Normalization Mathematics}
\author{Source: \texttt{repos/heterosplat/Documents/LaTeX/NormalizationMathematics.tex}}
\date{\today}

\newcommand{\SO}{\mathrm{SO}}
\newcommand{\R}{\mathbb{R}}
\newcommand{\HQ}{\mathbb{H}}
\newcommand{\diag}{\operatorname{diag}}

\begin{document}
\maketitle

\section{Purpose}

This document describes the mathematics behind the \textbf{Normalize} module
(Phase~2 of heterosplat). The module transforms 3D Gaussian splat parameters
from arbitrary coordinate conventions into a shared canonical frame, enabling
fusion of multiple splat captures.

\paragraph{Source files.}
\begin{itemize}
  \item \texttt{Source/Normalize/Transform.h} / \texttt{Transform.cu} ---
        CUDA kernels for similarity and homogeneous transforms.
  \item \texttt{Source/Normalize/Convention.h} / \texttt{Convention.cpp} ---
        up-axis detection, scene extent, normalization transform computation.
  \item \texttt{Source/NormalizeAndFuse.cu} --- binary that loads, normalizes,
        and optionally fuses two \texttt{.ply} files.
\end{itemize}

\section{Similarity transforms on Gaussians}
\label{sec:similarity}

A \emph{similarity transform} is parameterized as
\begin{equation}
  \mathbf{p}' = s \, R \, \mathbf{p} + \mathbf{t},
  \qquad R \in \SO(3),\quad s > 0,\quad \mathbf{t} \in \R^3.
  \label{eq:similarity}
\end{equation}
Each 3D Gaussian splat is defined by parameters
$(\boldsymbol{\mu},\, q,\, \boldsymbol{\sigma},\, \alpha,\, \mathbf{c})$:
mean position, orientation quaternion, log-scale vector, opacity, and
spherical harmonics coefficients. Under~\eqref{eq:similarity} each transforms
as follows.

\subsection{Means}
Direct application:
\begin{equation}
  \boldsymbol{\mu}' = s \, R \, \boldsymbol{\mu} + \mathbf{t}.
\end{equation}

\subsection{Quaternions (orientation)}
\label{sec:quat-compose}

Each Gaussian's 3D covariance is $\Sigma = R(q)\,S^2\,R(q)^T$ where
$R(q)$ is the rotation decoded from the Gaussian's quaternion~$q$ and
$S = \diag(e^{\sigma_1}, e^{\sigma_2}, e^{\sigma_3})$. The similarity
rotation acts by left-composing:
\begin{equation}
  \Sigma' = R\,R(q)\,S'^2\,R(q)^T\,R^T
  \quad\Longrightarrow\quad
  q' = q_R \otimes q,
\end{equation}
where $q_R \in \HQ_1$ is the unit quaternion corresponding to~$R$ and
$\otimes$ denotes Hamilton product.

\paragraph{Quaternion multiplication.}
Given $a = (a_w, a_x, a_y, a_z)$ and $b = (b_w, b_x, b_y, b_z)$, the
Hamilton product $c = a \otimes b$ is:
\begin{align}
  c_w &= a_w b_w - a_x b_x - a_y b_y - a_z b_z, \\
  c_x &= a_w b_x + a_x b_w + a_y b_z - a_z b_y, \\
  c_y &= a_w b_y - a_x b_z + a_y b_w + a_z b_x, \\
  c_z &= a_w b_z + a_x b_y - a_y b_x + a_z b_w.
\end{align}

\paragraph{Matrix-to-quaternion conversion (Shepperd's method).}
Given rotation matrix $R$ with entries $r_{ij}$, let
$\mathrm{tr} = r_{00} + r_{11} + r_{22}$. If $\mathrm{tr} > 0$:
\begin{equation}
  s = \frac{1}{2\sqrt{\mathrm{tr}+1}}, \quad
  q = \Bigl(\tfrac{1}{4s},\;
  (r_{21}-r_{12})s,\;
  (r_{02}-r_{20})s,\;
  (r_{10}-r_{01})s\Bigr).
\end{equation}
Otherwise, choose the diagonal element $r_{ii}$ with the largest value and
apply the corresponding formula (three cases for $i = 0, 1, 2$).

\subsection{Log-scales}

The uniform scale factor~$s$ multiplies each axis of the Gaussian equally:
\begin{equation}
  \boldsymbol{\sigma}' = \boldsymbol{\sigma} + \ln(s) \cdot \mathbf{1}_3.
\end{equation}
In code: \texttt{log\_scales[n*3+k] += logf(scale)} for $k = 0,1,2$.

\subsection{Opacity and SH coefficients}

Opacity is invariant under similarity transforms. The DC component of
spherical harmonics is also invariant (it represents view-independent radiance).

Higher-order SH coefficients ($\ell \geq 1$) should be rotated using Wigner
$D$-matrices of the rotation~$R$. For degree~$\ell$, the rotation acts on
the $(2\ell+1)$-dimensional subspace via
\begin{equation}
  c_\ell^{'m} = \sum_{m'=-\ell}^{\ell} D^{\ell}_{mm'}(R)\, c_\ell^{m'},
\end{equation}
where $D^\ell(R)$ is the real Wigner $D$-matrix. This is deferred in the
current implementation; the kernel only transforms means, quaternions, and
log-scales.

\section{Homogeneous transform (means only)}
\label{sec:homogeneous}

For non-rigid or non-uniform transforms where rotation extraction is ambiguous,
we apply a full $4 \times 4$ homogeneous matrix to means only:
\begin{equation}
  \begin{pmatrix} x' \\ y' \\ z' \\ w \end{pmatrix}
  = M \begin{pmatrix} x \\ y \\ z \\ 1 \end{pmatrix},
  \qquad
  \mathbf{p}' = \frac{1}{w}\begin{pmatrix} x' \\ y' \\ z' \end{pmatrix}.
\end{equation}
For affine transforms the last row is $(0,0,0,1)$ and $w = 1$. The
perspective division handles the general projective case.

\section{Coordinate convention normalization}
\label{sec:convention}

Different 3D capture tools use different conventions:
\begin{center}
\begin{tabular}{ll}
  \textbf{Tool} & \textbf{Convention} \\
  \hline
  COLMAP, engineering & Z-up \\
  Blender, OpenGL, Unity & Y-up \\
\end{tabular}
\end{center}

\subsection{Y-up to Z-up rotation}

The canonical target frame is Z-up. To convert from Y-up, rotate $-90°$
around the $X$-axis:
\begin{equation}
  R_{Y\to Z} =
  \begin{pmatrix}
    1 & 0 & 0 \\
    0 & 0 & -1 \\
    0 & 1 & 0
  \end{pmatrix}.
\end{equation}
This maps $(x, y, z)_{\text{Y-up}} \mapsto (x, -z, y)_{\text{Z-up}}$.

\subsection{Automatic detection}

The up-axis is detected heuristically from the point cloud statistics:
\begin{enumerate}
  \item Compute per-axis variance $\sigma^2_x, \sigma^2_y, \sigma^2_z$.
  \item The up-axis typically has the smallest variance (scenes are wider
        than tall).
  \item If the minimum-variance axis is clearly smaller ($< 50\%$ of the
        second-smallest), use it.
  \item Otherwise, fall back to checking which axis has a positive mean
        (objects are above the ground plane).
\end{enumerate}

\subsection{Unit-sphere normalization}

After up-axis alignment, the scene is centered and scaled to fit in a
unit sphere:
\begin{align}
  s &= \frac{2}{\max(\Delta x, \Delta y, \Delta z)}, \\
  \mathbf{t} &= -s \, R \, \bar{\mathbf{p}},
\end{align}
where $\bar{\mathbf{p}}$ is the centroid and $\Delta$ denotes extent
(max $-$ min) along each axis. This normalization ensures that two captures
at different physical scales and positions are brought into register.

\section{Fusion}

After both sources are normalized to the same frame (Z-up, unit sphere),
fusion is concatenation: the $N_A + N_B$ Gaussians are simply stacked.
Downstream rendering treats them as a single scene.

When the two sources have different SH degrees, the fused output uses the
minimum degree, truncating higher-order coefficients from the richer source.

\end{document}
